% ==============================================================================
% AR(2) extension of the baseline EM framework
% Created: 2026-05-05
%
% This document extends the baseline three-wave AR(1) EM algorithm
% (EM baseline.tex) to a four-wave AR(2) specification. True employment
% follows a second-order Markov chain; observed employment is measured
% with symmetric misclassification probability pi. The observation model,
% misclassification structure, and convergence properties are inherited
% from the baseline; only the latent dynamics, history space, E-step
% sufficient statistics, and M-step transition updates change.
%
% Notation is consistent with EM baseline.tex.
% ==============================================================================
\documentclass[12pt]{article}
\usepackage[margin=2cm]{geometry}
\setlength{\parindent}{1em}
\setlength{\parskip}{1em}

\usepackage{amsmath}
\usepackage{amssymb}
\usepackage{booktabs,dcolumn,setspace,float}

\begin{document}

\title{Second-order Markov extension of the EM algorithm for misclassified
employment transitions}
\author{}
\date{May 2026}
\maketitle

\section{Motivation}

The baseline specification models true employment as a first-order
Markov chain: the transition probability at wave~$t$ depends only on
$h_{t-1}$. This rules out duration dependence---the empirically
important phenomenon whereby the probability of exiting a state depends
on how long the individual has been in that state. In discrete time with
quarterly observations, the minimal extension that captures
one-period duration dependence is a second-order Markov chain, where
$\Pr(h_t \mid h_{t-1}, h_{t-2})$ depends on the two most recent states.
Testing for second-order dependence is equivalent to testing whether the
additional parameters are zero, providing a natural specification test
for the baseline model.

This section extends the baseline EM algorithm to the AR(2) case. The
observation model (symmetric misclassification with probability~$\pi$,
equation~\eqref{eq:misclass} of the baseline) is unchanged. The panel
is extended from three to four waves to provide sufficient transitions
for identification.

\section{Model}

\subsection{States and notation}

Consider a population observed at four equally spaced waves
$t \in \{1, 2, 3, 4\}$, separated by quarterly intervals of
$\Delta = 0.25$ years. As in the baseline, $h_t \in \{0,1\}$ denotes
the true (latent) employment state and $s_t \in \{0,1\}$ the observed
state.

The latent four-period history is
\[
H = (h_1, h_2, h_3, h_4) \in \{0,1\}^4,
\]
and the set of all $2^4 = 16$ possible histories is denoted
$\mathcal{H}$. The observed sequence is
$s = (s_1, s_2, s_3, s_4)$.

\subsection{True dynamics: second-order Markov chain}
\label{sec:ar2_dynamics}

The true employment process $\{h_t\}$ is a second-order Markov chain.
The transition probabilities are parameterised as
\begin{equation}
\Pr(h_t = 1 \mid h_{t-1}, h_{t-2}) =
\begin{cases}
\theta_0                          & \text{if } h_{t-2} = 0,\; h_{t-1} = 0, \\
\theta_0 + \theta_{01}            & \text{if } h_{t-2} = 1,\; h_{t-1} = 0, \\
1 - \theta_1 - \theta_{10}        & \text{if } h_{t-2} = 0,\; h_{t-1} = 1, \\
1 - \theta_1                      & \text{if } h_{t-2} = 1,\; h_{t-1} = 1,
\end{cases}
\label{eq:ar2_transitions}
\end{equation}
where $\theta_0 \in (0,1)$ is the baseline job-finding rate (from
two consecutive periods of nonemployment), $\theta_1 \in (0,1)$ is the
baseline separation rate (from two consecutive periods of employment),
and the \emph{duration-dependence increments} $\theta_{01}$ and
$\theta_{10}$ capture the effect of the state two periods ago. When
$\theta_{01} = \theta_{10} = 0$, the model reduces to the AR(1)
baseline. All four transition probabilities are constrained to lie in
$(0,1)$, which imposes $\theta_0 + \theta_{01} \in (0,1)$ and
$\theta_1 + \theta_{10} \in (0,1)$.

For notational convenience, define the four unrestricted transition
probabilities:
\begin{equation}
p_{00} \equiv \theta_0, \qquad
p_{10} \equiv \theta_0 + \theta_{01}, \qquad
p_{01} \equiv 1 - \theta_1 - \theta_{10}, \qquad
p_{11} \equiv 1 - \theta_1,
\label{eq:p_from_theta}
\end{equation}
where $p_{jk} = \Pr(h_t = 1 \mid h_{t-2} = j,\, h_{t-1} = k)$. The
inverse mapping recovers the structural parameters:
\begin{equation}
\theta_0 = p_{00}, \qquad
\theta_{01} = p_{10} - p_{00}, \qquad
\theta_1 = 1 - p_{11}, \qquad
\theta_{10} = p_{11} - p_{01}.
\label{eq:theta_from_p}
\end{equation}

\subsection{Stationary distribution}
\label{sec:stationary}

Under stationarity, the chain is in its ergodic steady state at waves~1
and~2. The second-order chain on $\{h_t\}$ is equivalent to a
first-order chain on the state pair $(h_{t-1}, h_t) \in \{0,1\}^2$,
with $4 \times 4$ transition matrix. The stationary distribution
$\alpha(h_1, h_2) \equiv \Pr(h_1, h_2)$ satisfies the balance equations.

Define the normalising constant
\begin{equation}
\Phi \equiv \theta_0(\theta_{10} - 1) + \theta_1(\theta_{01} - 1).
\label{eq:Phi}
\end{equation}
Then the joint stationary probabilities are
\begin{align}
\alpha(0, 0) &= \frac{\theta_1\,(\theta_0 + \theta_{01} - 1)}{\Phi},
\label{eq:mu00} \\[4pt]
\alpha(1, 0) &= \frac{-\theta_0\,\theta_1}{\Phi},
\label{eq:mu10} \\[4pt]
\alpha(0, 1) &= \frac{-\theta_0\,\theta_1}{\Phi},
\label{eq:mu01} \\[4pt]
\alpha(1, 1) &= \frac{\theta_0\,(\theta_1 + \theta_{10} - 1)}{\Phi}.
\label{eq:mu11}
\end{align}
Note that $\alpha(1,0) = \alpha(0,1)$, which follows from the detailed
balance of the embedded pair chain. The marginal employment rate at
any single wave is $\alpha(0,1) + \alpha(1,1)$.

\subsection{Prior probability of a latent history}

The prior probability of a history $h = (h_1, h_2, h_3, h_4)$ is
\begin{equation}
p(H = h \mid \Theta) =
\alpha(h_1, h_2)
\prod_{t=3}^{4}
p_{h_{t-2},\, h_{t-1}}^{\,\mathbf{1}\{h_t = 1\}}\,
(1 - p_{h_{t-2},\, h_{t-1}})^{\mathbf{1}\{h_t = 0\}},
\label{eq:ar2_prior}
\end{equation}
where $\alpha(h_1, h_2)$ is given by
\eqref{eq:mu00}--\eqref{eq:mu11} and $p_{jk}$ by
\eqref{eq:p_from_theta}. The product runs over waves~3 and~4, each
contributing one second-order transition factor.

\subsection{Observed-data likelihood}

The observation model is identical to the baseline: symmetric
misclassification with probability $\pi$ (equation~\eqref{eq:misclass}
of the baseline), independent across waves conditional on the latent
history. The marginal probability of the observed sequence
$y_i = (s_{i1}, s_{i2}, s_{i3}, s_{i4})$ is
\begin{equation}
p(y_i \mid \Theta) = \sum_{h \in \mathcal{H}}
  p(H = h \mid \Theta)\,
  \prod_{t=1}^{4} q_{s_{it},\, h_t}(\pi),
\label{eq:ar2_marginal}
\end{equation}
and the observed-data log-likelihood is
\begin{equation}
\ell(\Theta) = \sum_{i=1}^{N} w_i \log p(y_i \mid \Theta),
\label{eq:ar2_obs_ll}
\end{equation}
with parameter vector
$\Theta = (\theta_0, \theta_{01}, \theta_{10}, \theta_1, \pi)$.
This is a finite mixture over the 16 latent histories.

\section{EM algorithm}
\label{sec:ar2_em}

The EM algorithm follows the same structure as the baseline. The
complete-data log-likelihood decomposes additively into a
\emph{dynamics block} (depending on the transition parameters through
the AR(2) prior) and a \emph{misclassification block} (depending only
on~$\pi$), enabling separate maximisation in the M-step.

\subsection{E-step: responsibilities}

Starting from current estimates $\Theta^{(\mathrm{old})}$, the E-step
computes posterior responsibilities
\begin{equation}
\gamma_{ih} = \frac{
  p(H = h \mid \Theta^{(\mathrm{old})})\,
  \prod_{t=1}^{4} q_{s_{it},\, h_t}(\pi^{(\mathrm{old})})
}{
  \sum_{h' \in \mathcal{H}}
  p(H = h' \mid \Theta^{(\mathrm{old})})\,
  \prod_{t=1}^{4} q_{s_{it},\, h'_t}(\pi^{(\mathrm{old})})
},
\label{eq:ar2_responsibility}
\end{equation}
for each individual $i$ and history $h \in \mathcal{H}$.
Each numerator is a product of at most nine terms: four from the AR(2)
prior (one joint initial probability and two transition factors) and
four from misclassification. With $N$ individuals and 16 histories, the
E-step requires $O(N)$ operations.

\paragraph{Sufficient statistics.}
From the responsibilities, accumulate the following weighted counts.
For the transition parameters, define four origin-type counters and
four transition counters, indexed by the state pair
$(h_{t-2}, h_{t-1})$:
\begin{align}
D_{jk} &= \sum_{i=1}^{N} \sum_{h \in \mathcal{H}} w_i \gamma_{ih}
  \sum_{t=3}^{4} \mathbf{1}\{h_{t-2} = j,\, h_{t-1} = k\},
  \label{eq:ar2_suff_D} \\[4pt]
T_{jk} &= \sum_{i=1}^{N} \sum_{h \in \mathcal{H}} w_i \gamma_{ih}
  \sum_{t=3}^{4} \mathbf{1}\{h_{t-2} = j,\, h_{t-1} = k,\, h_t = 1\},
  \label{eq:ar2_suff_T}
\end{align}
for $j, k \in \{0, 1\}$. Here $D_{jk}$ counts
(responsibility-weighted) person-periods where the two preceding states
were $(j, k)$, and $T_{jk}$ counts the subset of those that transition
to employment.

For the misclassification parameter:
\begin{equation}
M = \sum_{i=1}^{N} \sum_{h \in \mathcal{H}} w_i \gamma_{ih}
  \sum_{t=1}^{4} \mathbf{1}\{s_{it} \neq h_t\}.
  \label{eq:ar2_suff_M}
\end{equation}

\subsection{M-step: parameter updates}
\label{sec:ar2_mstep}

The M-step maximises $Q(\Theta \mid \Theta^{(\mathrm{old})})$ by
updating each block independently.

\paragraph{Transition probabilities (closed form).}
As in the baseline, we adopt the standard approximation for short
panels: update the transition parameters from the transition counts
alone, then set the initial distribution by stationarity. The
unrestricted transition probabilities have closed-form updates:
\begin{equation}
\hat{p}_{jk} = \frac{T_{jk}}{D_{jk}}, \qquad j, k \in \{0, 1\}.
\label{eq:ar2_mstep_p}
\end{equation}
Each $\hat{p}_{jk}$ is the responsibility-weighted fraction of
person-periods with origin pair $(j, k)$ that result in employment.
The structural parameters are then recovered via
\eqref{eq:theta_from_p}:
\begin{equation}
\hat{\theta}_0 = \hat{p}_{00}, \qquad
\hat{\theta}_{01} = \hat{p}_{10} - \hat{p}_{00}, \qquad
\hat{\theta}_1 = 1 - \hat{p}_{11}, \qquad
\hat{\theta}_{10} = \hat{p}_{11} - \hat{p}_{01}.
\label{eq:ar2_mstep_theta}
\end{equation}
The stationary distribution $\alpha(h_1, h_2)$ is updated from the new
$\hat{\theta}$ values via
\eqref{eq:Phi}--\eqref{eq:mu11}.

\paragraph{Misclassification probability (closed form).}
\begin{equation}
\hat{\pi} = \frac{M}{4\,W},
\label{eq:ar2_mstep_pi}
\end{equation}
where $W = \sum_{i=1}^{N} w_i$. The denominator is $4W$ (rather than
$3W$ in the baseline) because each individual now contributes four
wave-observations.

\paragraph{Summary.} The complete M-step consists of five closed-form
updates: four transition probabilities
\eqref{eq:ar2_mstep_p} and one misclassification probability
\eqref{eq:ar2_mstep_pi}. No numerical optimisation is required.

\subsection{Iteration and convergence}

The iteration scheme and convergence properties are identical to the
baseline (Section~\ref{sec:convergence} of the baseline). The EM
guarantees monotone ascent of $\ell(\Theta)$; all M-step updates are in
closed form; the observed-data likelihood is bounded above; and every
limit point of the parameter sequence is a stationary point. The same
practical recommendations apply: multiple starts to guard against local
maxima, boundary constraints
$p_{jk} \in [\varepsilon, 1 - \varepsilon]$ and
$\pi \in [\varepsilon, \tfrac{1}{2} - \varepsilon]$, and survey weights
throughout.

\subsection{Testing for second-order dependence}

The AR(1) model is nested within the AR(2) model under the restriction
$\theta_{01} = \theta_{10} = 0$, or equivalently
$p_{10} = p_{00}$ and $p_{01} = p_{11}$. This can be tested via a
likelihood ratio test:
\begin{equation}
\mathrm{LR} = 2\big[\ell(\hat\Theta_{\mathrm{AR(2)}})
  - \ell(\hat\Theta_{\mathrm{AR(1)}})\big]
\;\sim\; \chi^2_2 \quad \text{under } H_0: \theta_{01} = \theta_{10} = 0.
\label{eq:ar2_lr_test}
\end{equation}

\section{Identification}
\label{sec:ar2_identification}

The model has five free parameters,
$\Theta = (\theta_0, \theta_{01}, \theta_{10}, \theta_1, \pi)$.
The data provide $2^4 - 1 = 15$ free moments (the probabilities of the
16 observed patterns $(s_1, s_2, s_3, s_4)$, summing to one). The model
is therefore overidentified with ten overidentifying restrictions.

\paragraph{Identification without misclassification.}
When $\pi = 0$, observed states equal true states, and the four
transition probabilities $p_{jk}$ are directly identified from the
observed second-order transition frequencies:
\[
p_{jk} = \Pr(s_t = 1 \mid s_{t-2} = j,\, s_{t-1} = k),
\]
which can be estimated from the eight second-order transition counts
across waves~2--3 and~3--4. With 15 free moments and four parameters,
the model without misclassification has 11 overidentifying restrictions.

\paragraph{Identification with misclassification.}
When $\pi > 0$, the key identifying argument parallels the baseline.
In a true second-order Markov chain, $h_4$ is conditionally independent
of $h_1$ given $(h_2, h_3)$. Misclassification breaks this conditional
independence in the observed data: conditioning on $(s_2, s_3)$ does not
fully screen off $s_1$ from $s_4$ because each $s_t$ is a noisy signal
of~$h_t$. The residual dependence between $s_1$ and $s_4$ given
$(s_2, s_3)$ provides identifying information for~$\pi$.

More formally, the model maps five parameters to 15 free moments.
The Jacobian of the map
$\Theta \mapsto \{p_{abcd}(\Theta)\}$ has full column rank~5
throughout the interior of the parameter space (verified analytically
via symbolic computation), confirming generic identification. The ten
overidentifying restrictions can be tested via a likelihood ratio test
against the saturated model:
\begin{equation}
\mathrm{LR} = 2\big[\ell_{\mathrm{saturated}} - \ell(\hat\Theta)\big]
\;\sim\; \chi^2_{10} \quad \text{under } H_0.
\label{eq:ar2_overid_test}
\end{equation}
Rejection indicates that the second-order Markov assumption with
symmetric misclassification does not adequately describe the data,
motivating further extensions such as the tenure-augmented
contamination model.

\end{document}
